\section{Background: Speech LLMs and Their Lifecycle Challenges}
\label{sec:background}

We assume reader familiarity with the Transformer-based language model
backbone (e.g., \citep{vaswani2017attention,radford2019language,brown2020language})
and standard supervised, instruction-tuned, and RLHF-aligned text LLM
pipelines. This section sets up the speech-specific concepts the rest of
the survey relies on.

\paragraph{Speech LLMs.}
We use \emph{Speech LLM} to refer to an LLM-based system that consumes,
reasons over, or generates speech as part of language interaction. This
includes speech-to-text recognition coupled with an LLM (cascaded),
speech-to-text and speech-to-speech models that share a single backbone
(end-to-end), and omni-modal models that integrate speech with text and
other modalities \citep{cui2024recent,peng2025surveyspeechlargelanguage}.

\paragraph{Why speech is different.}
Three properties of speech shape every lifecycle stage of Speech LLMs.
First, speech is a continuous, high-dimensional, multi-channel signal that
carries linguistic content, paralinguistics, speaker identity, scene
context, and timing simultaneously. Second, parallel speech-text data is
expensive and unevenly distributed across languages and domains, so most
training pipelines depend on weak supervision or self-supervision
\citep{radford2023robust,hsu2021hubert}. Third, speech interaction is
inherently real-time: a model that produces a correct transcript or
response too late may fail the interaction even if its content is
accurate. Together, these properties drive the design choices reviewed in
the following sections.

\paragraph{Lifecycle framing.}
The remainder of the survey is organized around four lifecycle stages:
training strategies (Section~\ref{sec:training_strategies}), evaluation
and benchmarking (Section~\ref{sec:evaluation_benchmarking}), deployment
and real-time interaction (Section~\ref{sec:deployment}), and safety,
privacy, and governance (Section~\ref{sec:safety_privacy_governance}).
The four stages are tightly coupled. Training choices determine what
benchmarks reveal; benchmarks reveal what deployment must defend
against; deployment constraints reshape which training and evaluation
trade-offs are acceptable; and safety risks span all three.
